\paragraph{Phương trình sóng tại điểm M cách nguồn O một khoảng x}
\begin{itemize}
\item Phương trình sóng tại nguồn O viết dưới dạng tổng quát: $u_O=A\cos(\omega t+\varphi)$.
\item Điểm M cách O một khoảng x có phương trình sóng dạng:
\begin{align*}
\boder{u_M =A\cos(\omega t+\varphi - \omega\dfrac{x}{v})= A\cos(\omega t+\varphi - \dfrac{2\pi x}{\lambda})};\ t\geq \dfrac{x}{v}
\end{align*}
\end{itemize}
\paragraph{Phương trình sóng tại điểm Q cách điểm P cho trước một khoảng x trên phương truyền sóng}
\Binhluan{\begin{itemize}
\item Phương trình sóng tại P viết dưới dạng tổng quát: \begin{align*}u_P=A\cos(\omega t+\varphi)\end{align*}
\item Điểm Q cách P một khoảng x có phương trình sóng dạng: \begin{align*}\boder{u_Q=A\cos(\omega t+\varphi \pm \dfrac{2\pi x}{\lambda})}\end{align*}
\end{itemize}}{Xét theo chiều truyền sóng, nếu Q trước P thì lấy $+\dfrac{2\pi x}{\lambda}$; nếu Q sau P thì lấy $-\dfrac{2\pi x}{\lambda}$}
\newghichu{
Hàm sóng là một hàm vừa tuần hoàn theo thời gian với chu kì T, vừa tuần hoàn theo không gian với khoảng cách $\lambda$.
}
